\section{Results}\label{sec:results}

\subsection{Data readiness and label frame}

The v8 integration completed without row loss: all 16,506 grid units are
present and 374 fields are available. Remote-sensing, surface-model, and climate-response
families have low missingness, whereas building, floor, frontage, and network
families show the expected structural absence where the relevant
objects. The highest family-average missing share is 36.3\% for floor proxies;
surface-model missingness averages 0.7\% and canopy height 1.8\%. Figure~\ref{fig:data-readiness}
contrasts feature-family missingness with the four-class analysis population.
composition. LCZ 6 is the largest class (1,842), followed by LCZ 8 (1,317), LCZ
9 (1,099), and LCZ 3 (1,081).

The manual audit is complete: all 569 sampled units passed the schema gate, comprising
382 high-, 143 medium-, and 28 low-quality rows, with 16 edge-case reject flags.
Weak-versus-audit agreement and bounded audited model comparisons are therefore
available, although the primary four-class comparison is limited to 96
high-quality units and retains substantial class-space compression.

\begin{figure}[H]
    \centering
    \includegraphics[width=\textwidth]{fig_3_data_readiness.png}
    \caption{Data readiness and audit design. Panels (a--c) show representative
    250 m grid, LST, imagery, and feature evidence; (d) maps the complete
    569-unit audit sample by review pass; and (e) compares audit-quality flags
    with priority scores. Missingness can be structural and is not automatically
    an error.}
    \label{fig:data-readiness}
\end{figure}

\subsection{Weak-class climate-morphology profiles}

The retained weak classes were morphologically distinct (Table~\ref{tab:class_profiles};
Figure~\ref{fig:class-profiles}). LCZ 3 had the highest median building coverage
(0.380) and floor-height proxy (11.10 m). LCZ 8 combined intermediate coverage
(0.191) and height (8.33 m) with the highest summer median LST, 46.09~$^{\circ}$C.
LCZ 6 and LCZ 9 were much more open: median coverage was 0.024 and 0.001,
respectively, and median LST was 43.01 and 43.27~$^{\circ}$C. LCZ 6 and LCZ 9
also had higher median NDVI (0.331 and 0.353) than LCZ 3 and LCZ 8 (0.137 and
0.132).

\begin{table*}[htbp]
\centering
\footnotesize
\caption{Median climate-morphology profiles of the four retained weak-label classes.}
\label{tab:class_profiles}
\resizebox{\textwidth}{!}{%
\begin{tabular}{lrrrrrrrr}
\toprule
Class & n & LST & Coverage (\%) & Height proxy (m) & Surface-model SD (m) & Canopy (\%) & NDVI & 2SFCA log1p \\
\midrule
LCZ 3 & 1,081 & 44.85 & 38.0 & 11.10 & 7.26 & 100.0 & 0.137 & 0.567 \\
LCZ 6 & 1,842 & 43.01 & 2.4 & 6.37 & 9.30 & 91.3 & 0.331 & 0.000 \\
LCZ 8 & 1,317 & 46.09 & 19.1 & 8.33 & 5.17 & 87.7 & 0.132 & 0.217 \\
LCZ 9 & 1,099 & 43.27 & 0.1 & 3.20 & 9.87 & 80.7 & 0.353 & 0.000 \\
\bottomrule
\end{tabular}%
}
\begin{minipage}{0.96\textwidth}\vspace{2pt}\scriptsize Note: Descriptive weak-label profiles; they are not manually audited local subtypes.\end{minipage}
\end{table*}



\begin{figure}[H]
    \centering
    \includegraphics[width=0.78\textwidth]{fig_4_class_profiles.png}
    \caption{Weak-class profiles. Colour expresses each class median standardized
    across the four classes; annotations are raw medians. Structural missingness
means the displayed height-proxy medians use available building-bearing units.}
    \label{fig:class-profiles}
\end{figure}

These contrasts are consistent with the broad semantics of the global labels,
but they remain weak-label profiles rather than a substitute for the audited
local class space. In particular, canopy
cover above 2 m is high in all four classes because the ETH surface-height
product and threshold capture vegetation pixels wherever present; it should not
be read as public green-space share. Likewise, road density is a geometric
network measure rather than walkability.

\subsection{Spatial-block baseline performance and ablation}

The majority classifier produced macro-F1 0.128. LightGBM morphology-only
macro-F1 was 0.663; adding vegetation raised it to 0.760, and green-blue
context raised it to 0.802. The full recipe without 2SFCA reached 0.833, while
the preregistered 800 m recipe reached 0.838 for LightGBM, 0.835 for XGBoost,
0.829 for Random Forest, and 0.819 for Extra Trees (Figure~\ref{fig:baseline-performance};
Table~\ref{tab:baseline_performance}). The LightGBM standard deviation at the
primary recipe was 0.012. The 1,200 m alternative was marginally higher for
LightGBM (0.839), so the 800 m result is retained as the prespecified primary
rather than selected post hoc.

\begin{figure}[!htbp]
    \centering
    \includegraphics[width=0.92\textwidth]{fig_5_baseline_performance.png}
    \caption{Mean macro-F1 under five spatial block folds for four tabular
    learners. Recipes are additive and the shaded position marks the
    prespecified 800 m recipe. All targets are global weak labels.}
    \label{fig:baseline-performance}
\end{figure}

\begin{table*}[htbp]
\centering
\scriptsize
\caption{Five-fold stratified spatial-block weak-label performance.}
\label{tab:baseline_performance}
\setlength{\tabcolsep}{3pt}
\resizebox{\textwidth}{!}{%
\begin{tabular}{llrrr}
\toprule
Model & Recipe & Features & Macro-F1 (mean $\pm$ SD) & Balanced accuracy \\
\midrule
XGBoost & Morphology & 15 & 0.662 $\pm$ 0.037 & 0.678 \\
LightGBM & Morphology & 15 & 0.663 $\pm$ 0.027 & 0.669 \\
Extra Trees & Morphology & 15 & 0.673 $\pm$ 0.030 & 0.705 \\
Random Forest & Morphology & 15 & 0.676 $\pm$ 0.029 & 0.688 \\
LightGBM & + vegetation & 22 & 0.760 $\pm$ 0.010 & 0.762 \\
XGBoost & + vegetation & 22 & 0.761 $\pm$ 0.023 & 0.766 \\
Extra Trees & + vegetation & 22 & 0.762 $\pm$ 0.016 & 0.776 \\
Random Forest & + vegetation & 22 & 0.767 $\pm$ 0.015 & 0.772 \\
Extra Trees & + green-blue & 29 & 0.795 $\pm$ 0.009 & 0.803 \\
XGBoost & + green-blue & 29 & 0.801 $\pm$ 0.013 & 0.805 \\
LightGBM & + green-blue & 29 & 0.802 $\pm$ 0.011 & 0.804 \\
Random Forest & + green-blue & 29 & 0.805 $\pm$ 0.010 & 0.808 \\
Extra Trees & Full, no 2SFCA & 67 & 0.819 $\pm$ 0.015 & 0.824 \\
Random Forest & Full, no 2SFCA & 67 & 0.827 $\pm$ 0.011 & 0.827 \\
LightGBM & Full, no 2SFCA & 67 & 0.833 $\pm$ 0.013 & 0.834 \\
XGBoost & Full, no 2SFCA & 67 & 0.833 $\pm$ 0.013 & 0.836 \\
Extra Trees & Full + 400 m & 68 & 0.819 $\pm$ 0.017 & 0.824 \\
Random Forest & Full + 400 m & 68 & 0.824 $\pm$ 0.012 & 0.825 \\
LightGBM & Full + 400 m & 68 & 0.831 $\pm$ 0.012 & 0.832 \\
XGBoost & Full + 400 m & 68 & 0.835 $\pm$ 0.015 & 0.837 \\
Extra Trees & Full + 800 m & 68 & 0.819 $\pm$ 0.017 & 0.825 \\
Random Forest & Full + 800 m & 68 & 0.829 $\pm$ 0.012 & 0.830 \\
XGBoost & Full + 800 m & 68 & 0.835 $\pm$ 0.010 & 0.838 \\
LightGBM & Full + 800 m & 68 & 0.838 $\pm$ 0.012 & 0.839 \\
Extra Trees & Full + 1200 m & 68 & 0.821 $\pm$ 0.016 & 0.826 \\
Random Forest & Full + 1200 m & 68 & 0.828 $\pm$ 0.013 & 0.828 \\
XGBoost & Full + 1200 m & 68 & 0.838 $\pm$ 0.013 & 0.841 \\
LightGBM & Full + 1200 m & 68 & 0.839 $\pm$ 0.013 & 0.840 \\
\bottomrule
\end{tabular}%
}
\begin{minipage}{0.96\textwidth}\vspace{2pt}\scriptsize Note: Blocks are nominally 1.25 km (5 by 5 cells). Preprocessing is fit within each training fold. Values are mean $\pm$ standard deviation across five folds. These are exploratory weak-label diagnostics; within each recipe, models are ordered by descending macro-F1.\end{minipage}
\end{table*}


Pooled out-of-fold LightGBM F1 was 0.861 for LCZ 3, 0.833 for LCZ 6, 0.820 for
LCZ 8, and 0.839 for LCZ 9 (Table~\ref{tab:class_metrics}). The lower F1 for
LCZ 8 indicates overlap or heterogeneity around large low-rise fabrics, even
though the class has the highest LST profile.

\begin{table}[!htbp]
\centering
\scriptsize
\caption{Pooled out-of-fold class performance for the primary 800 m recipe.}
\label{tab:class_metrics}
\resizebox{\textwidth}{!}{%
\begin{tabular}{llrrrr}
\toprule
Model & Class & Support & Precision & Recall & F1 \\
\midrule
LightGBM & LCZ 3 & 1081 & 0.852 & 0.870 & 0.861 \\
LightGBM & LCZ 6 & 1842 & 0.833 & 0.834 & 0.833 \\
LightGBM & LCZ 8 & 1317 & 0.838 & 0.803 & 0.820 \\
LightGBM & LCZ 9 & 1099 & 0.828 & 0.850 & 0.839 \\
XGBoost & LCZ 3 & 1081 & 0.850 & 0.873 & 0.862 \\
XGBoost & LCZ 6 & 1842 & 0.837 & 0.819 & 0.828 \\
XGBoost & LCZ 8 & 1317 & 0.828 & 0.800 & 0.813 \\
XGBoost & LCZ 9 & 1099 & 0.817 & 0.860 & 0.838 \\
Random Forest & LCZ 3 & 1081 & 0.843 & 0.886 & 0.864 \\
Random Forest & LCZ 6 & 1842 & 0.813 & 0.828 & 0.821 \\
Random Forest & LCZ 8 & 1317 & 0.840 & 0.778 & 0.808 \\
Random Forest & LCZ 9 & 1099 & 0.823 & 0.829 & 0.826 \\
Extra Trees & LCZ 3 & 1081 & 0.819 & 0.896 & 0.856 \\
Extra Trees & LCZ 6 & 1842 & 0.822 & 0.792 & 0.807 \\
Extra Trees & LCZ 8 & 1317 & 0.834 & 0.759 & 0.795 \\
Extra Trees & LCZ 9 & 1099 & 0.792 & 0.851 & 0.820 \\
\bottomrule
\end{tabular}%
}
\end{table}


\subsection{Feature importance and out-of-fold uncertainty}

LightGBM gain importance was distributed across remote semantic, 2D form,
vegetation, climate context, 3D/surface, street-network, and green-access
families (Figure~\ref{fig:importance}). Remote semantics contributed 48.0\% of
the normalized primary-recipe importance, followed by 2D form (20.2\%),
vegetation (10.4\%), and context (8.9\%); 2SFCA contributed 1.3\%. The leading
individual variables included WorldCover built-up share, building perimeter
density, Sentinel-2 NDWI, coast distance, Dynamic World crop probability, and
NDVI. This mixture supports multi-family integration, but it also confirms that
remote semantic information may share signal with the weak LCZ target. The
importance ranking is therefore an explanation of the model, not proof of an
independent physical mechanism.

\begin{figure}[!htbp]
    \centering
    \includegraphics[width=0.88\textwidth]{fig_6_feature_importance.png}
    \caption{Mean LightGBM gain importance across the five primary 800 m folds,
    normalized within each fold. Error bars are between-fold standard deviations.}
    \label{fig:importance}
\end{figure}

The primary LightGBM mean out-of-fold maximum probability was 0.959, compared
with pooled accuracy 0.837; correct predictions numbered 4,468 and mismatches
871. Figure~\ref{fig:uncertainty} reveals spatial clusters of lower confidence
and audit-compatible/contradictory units. Probability is not calibrated
confidence, however; the warning score is used as a stability indicator rather
than a planning-risk estimate.

\begin{figure}[!htbp]
    \centering
    \includegraphics[width=\textwidth]{fig_7_uncertainty.png}
    \caption{Out-of-fold LightGBM uncertainty for the primary recipe. (a) Spatial
    warning score with the completed high-quality audit overlay. (b) Weak-label,
    LightGBM, and Random Forest audit crosswalks. Agreement shares are
    descriptive and the probabilities are not calibrated confidence.}
    \label{fig:uncertainty}
\end{figure}

\subsection{Network 2SFCA threshold sensitivity}

Compared with the same full model without 2SFCA, the primary 800 m field
increased mean macro-F1 by 0.0056 for LightGBM, 0.0023 for XGBoost, 0.0028 for
Random Forest, and 0.0007 for Extra Trees. The 1,200 m alternative was slightly
higher for LightGBM (0.0064 above no 2SFCA), while threshold effects varied by
model. Figure~\ref{fig:2sfca-sensitivity} shows why these thresholds are treated
as sensitivity parameters: access geography changes with catchment radius, and
larger radii add both supply and competing demand. The 800 m threshold remains
primary by design rather than by post-hoc reselection.

\begin{figure}[!htbp]
    \centering
    \includegraphics[width=\textwidth]{fig_8_2sfca_sensitivity.png}
    \caption{Network-cost 2SFCA scale sensitivity. (a) Access geography at
    400, 800, and 1,200 m. (b) The 1,200 m minus 400 m access difference. (c)
    Zoomed examples of strongest gain, strongest 400 m advantage, and local
    contestation. The figure is descriptive of access geography; model-score
    deltas are reported in Table~\ref{tab:2sfca}.}
    \label{fig:2sfca-sensitivity}
\end{figure}

\begin{table*}[htbp]
\centering
\footnotesize
\caption{Paired 2SFCA threshold contribution relative to the same model without 2SFCA.}
\label{tab:2sfca}
\resizebox{\textwidth}{!}{%
\begin{tabular}{llrrrr}
\toprule
Model & Threshold & Mean delta & SD & Fold wins & Range \\
\midrule
LightGBM & 400 m & -0.0014 & 0.0033 & 1/5 & [-0.0057, +0.0036] \\
LightGBM & 800 m & +0.0056 & 0.0036 & 5/5 & [+0.0017, +0.0088] \\
LightGBM & 1200 m & +0.0064 & 0.0079 & 3/5 & [-0.0021, +0.0171] \\
XGBoost & 400 m & +0.0017 & 0.0080 & 3/5 & [-0.0059, +0.0129] \\
XGBoost & 800 m & +0.0023 & 0.0067 & 3/5 & [-0.0053, +0.0106] \\
XGBoost & 1200 m & +0.0054 & 0.0061 & 4/5 & [-0.0043, +0.0110] \\
Random Forest & 400 m & -0.0022 & 0.0052 & 2/5 & [-0.0095, +0.0028] \\
Random Forest & 800 m & +0.0028 & 0.0035 & 4/5 & [-0.0030, +0.0054] \\
Random Forest & 1200 m & +0.0014 & 0.0030 & 3/5 & [-0.0021, +0.0050] \\
Extra Trees & 400 m & +0.0003 & 0.0019 & 4/5 & [-0.0028, +0.0021] \\
Extra Trees & 800 m & +0.0007 & 0.0040 & 3/5 & [-0.0058, +0.0035] \\
Extra Trees & 1200 m & +0.0021 & 0.0060 & 3/5 & [-0.0047, +0.0094] \\
\bottomrule
\end{tabular}%
}
\begin{minipage}{0.96\textwidth}\vspace{2pt}\scriptsize Note: Threshold performance is interpreted as sensitivity, not post-hoc model selection; the 800 m threshold is primary by design. Deltas are paired within the same spatial folds.\end{minipage}
\end{table*}


The municipal source audit strengthens only the macro spatial interpretation.
The official dataset contained 189 distinct park sites; 180 sites in 112 matched
FUR neighbourhoods were comparable with 644 OSM polygons across 12 districts.
Pearson correlations on log counts and areas were 0.931 and 0.910, and
leave-one-district-out ranges were 0.882--0.945 and 0.808--0.928. Rank
correlations were lower (0.747 and 0.776). Only 10.6\% of comparable official
sites had exact or strong OSM name candidates, and only 417 of 1,237 total OSM
supply polygons were named. OSM supply therefore reproduces broad district
geography but cannot be called complete or equivalent to official public green
space.

\subsection{Leave-district-out transfer}

LODO LightGBM macro-F1 averaged 0.723 across 11 districts, with a standard
deviation of 0.104 and range 0.475--0.819 (Figure~\ref{fig:lodo};
Table~\ref{tab:leave_district}). G\"uzelbah\c{c}e (0.475), Konak (0.598), and
Bal\c{c}ova (0.682) were the hardest transfers, while Kar\c{s}\i yaka (0.819),
Karaba\u{g}lar (0.807), Narl\i dere (0.794), and Bayrakl\i{} (0.775) were
stronger. Mean accuracy was 0.818 and mean balanced accuracy 0.766. The decline
from spatial-block macro-F1 0.838 to district-held-out 0.723 indicates that
within-FUR geographic transfer is materially harder than local block
interpolation.

\begin{figure}[!htbp]
    \centering
    \includegraphics[width=0.9\textwidth]{fig_10_leave_district_out.png}
    \caption{LightGBM leave-one-district-out transfer. Only districts with at
    least 100 eligible observations are included; scores use classes present in each
    held-out district, with Random Forest shown as a comparison in panel (b).}
    \label{fig:lodo}
\end{figure}

\begin{table*}[htbp]
\centering
\footnotesize
\caption{Leave-one-district-out LightGBM weak-label robustness.}
\label{tab:leave_district}
\resizebox{\textwidth}{!}{%
\begin{tabular}{lrrrrr}
\toprule
Held-out district & $n$ & Classes & Accuracy & Balanced accuracy & Macro-F1 \\
\midrule
G\"uzelbah\c{c}e & 334 & 3 & 0.746 & 0.606 & 0.475 \\
Konak & 266 & 3 & 0.831 & 0.817 & 0.598 \\
Bal\c{c}ova & 161 & 4 & 0.795 & 0.684 & 0.682 \\
Cig\u{g}li & 367 & 4 & 0.823 & 0.779 & 0.723 \\
Gaziemir & 438 & 4 & 0.788 & 0.735 & 0.738 \\
Bornova & 1419 & 4 & 0.780 & 0.764 & 0.771 \\
Buca & 1082 & 4 & 0.807 & 0.783 & 0.774 \\
Bayrakli & 322 & 4 & 0.792 & 0.768 & 0.775 \\
Narl\i dere & 169 & 3 & 0.947 & 0.866 & 0.794 \\
Karaba\u{g}lar & 390 & 4 & 0.851 & 0.799 & 0.807 \\
Kar\c{s}\i yaka & 331 & 4 & 0.843 & 0.822 & 0.819 \\
\bottomrule
\end{tabular}%
}
\begin{minipage}{0.96\textwidth}\vspace{2pt}\scriptsize Note: Only districts with at least 100 eligible observations were evaluated. Macro-F1 uses classes present in the held-out district; LightGBM mean 0.723 and range 0.475--0.819. Random Forest, evaluated identically as a comparison baseline, reached mean macro-F1 0.700.\end{minipage}
\end{table*}


\subsection{Climate-response validation}

Summer LST differed across the four weak classes ($H=1592.0$, $p<0.001$,
$\epsilon^2=0.298$, $n=5,337$). LCZ 8 was warmest by median, followed by LCZ 3,
LCZ 9, and LCZ 6. Across grid units, NDBI had the strongest positive bivariate
association with LST ($\rho=0.684$), followed by building coverage
($\rho=0.442$), road density ($\rho=0.374$), and endpoint-derived intersection
density ($\rho=0.342$). NDVI ($\rho=-0.678$) and canopy-volume proxy
($\rho=-0.670$) had the strongest negative associations. Figure~\ref{fig:climate}
and Table~\ref{tab:climate_validation} report the complete selected set.

\begin{figure}[!htbp]
    \centering
    \includegraphics[width=\textwidth]{fig_9_climate_validation.png}
    \caption{Descriptive climate-response validation. (a) Summer median LST by
    weak LCZ class. (b) Spearman correlations between selected features and LST;
    the lower tiles show audited cool/hot contrasts. The statistics are
    cross-sectional and do not estimate causal cooling.}
    \label{fig:climate}
\end{figure}

\begin{table*}[htbp]
\centering
\footnotesize
\caption{Descriptive climate-response associations within the four-class weak-label subset.}
\label{tab:climate_validation}
\resizebox{\textwidth}{!}{%
\begin{tabular}{lrrr}
\toprule
Feature & n & Spearman rho with LST & p \\
\midrule
s2 ndbi mean & 5336 & +0.684 & <0.001 \\
s2 ndvi mean & 5336 & -0.678 & <0.001 \\
canopy volume gt2m proxy m3 per ha & 5337 & -0.670 & <0.001 \\
building coverage exact & 5337 & +0.442 & <0.001 \\
road density exact m per km2 & 5337 & +0.374 & <0.001 \\
network intersection density per km2 & 5337 & +0.342 & <0.001 \\
surface-model elevation m std & 5337 & -0.314 & <0.001 \\
canopy cover gt2m share & 5273 & -0.251 & <0.001 \\
morph open space fragmentation index & 5337 & +0.192 & <0.001 \\
green 2sfca 800m access log1p & 5056 & +0.178 & <0.001 \\
\bottomrule
\end{tabular}%
}
\begin{minipage}{0.96\textwidth}\vspace{2pt}\scriptsize Note: Kruskal-Wallis across LCZ 3/6/8/9: H=1592.0, p<0.001, epsilon-squared=0.298. Associations are bivariate and non-causal.\end{minipage}
\end{table*}



\subsection{Audited-label comparison}

The completed audit exposes a substantial mismatch between the four-class weak
label space and the audited local fabric codebook. In the prespecified
high-quality tier, 96 audited units were available: exact agreement was 8.3\% for the
weak label, 8.3\% for LightGBM, and 9.4\% for Random Forest. The high-plus-medium
sensitivity tier contained 148 audited units, with corresponding agreement of
10.1\%, 12.2\%, and 12.2\%; the all-quality four-class reference contained 165
units, with 10.3\%, 13.3\%, and 13.3\%. Figure~\ref{fig:audited-validation} shows
that many audited natural/open classes are not representable by the four-class
target. These values are therefore evidence of class-space compression and
model transfer limits, not a calibrated estimate of local subtype prevalence.

\begin{figure}[!htbp]
    \centering
    \includegraphics[width=\textwidth]{fig_11_audited_validation.png}
    \caption{Audited-label comparison. (a) Weak label versus audited local
    fabric, (b) LightGBM prediction versus audit, and (c) exact agreement by
    audit-quality tier. The primary tier is the 96-unit high-quality subset;
    agreement is descriptive and natural/open audited classes are retained to
    reveal class-space compression.}
    \label{fig:audited-validation}
\end{figure}

\begin{table*}[htbp]
\centering
\footnotesize
\caption{Manual-audit exact-class agreement for the four-class modeling population.}
\label{tab:audited_validation}
\resizebox{\textwidth}{!}{%
\begin{tabular}{lrrrr}
\toprule
Audit-quality tier & $n$ & Weak label (\%) & LightGBM (\%) & Random Forest (\%) \\
\midrule
High quality (primary) & 96 & 8.3 & 8.3 & 9.4 \\
High + medium (sensitivity) & 148 & 10.1 & 12.2 & 12.2 \\
All quality (reference) & 165 & 10.3 & 13.3 & 13.3 \\
\bottomrule
\end{tabular}%
}
\begin{minipage}{0.96\textwidth}\vspace{2pt}\scriptsize Note: The four-class modeling population is restricted to weak LCZ 3/6/8/9 with confidence $\geq 0.60$. Tiers were prespecified before agreement was computed. The broader unrestricted 569-unit audit had 36.9\% exact weak-label agreement and is not directly comparable to the built-subtype task.\end{minipage}
\end{table*}


\subsection{Claim-status synthesis}

Table~\ref{tab:evidence_status} separates supported, exploratory, and untestable
claims. The study currently supports reproducible data integration,
weak-label discrimination, network-access sensitivity, transfer heterogeneity,
descriptive climate relevance, and bounded audited comparison. It does not
support accurate audited local subtype mapping, deep multimodal superiority, or
planning-ready subtype prescriptions. The low audit agreement is itself a
result, not a basis for filling the remaining claims with extrapolation.

\begin{table}[!htbp]
\centering
\footnotesize
\caption{Analytical evidence summary.}
\label{tab:evidence_status}
\resizebox{\textwidth}{!}{%
\begin{tabular}{lll}
\toprule
Analytical component & Available evidence & Interpretation \\
\midrule
250 m data integration & 16,506 rows, 374 fields, no join loss & supported \\
Surface and context contribution & Surface summaries, ablations, and importance & interpretable contribution \\
Green-space accessibility & network-cost 2SFCA + reference macro audit & proxy supported \\
Spatial robustness & 5 folds + 11 held-out districts & weak-label supported \\
Climate relevance & LST class contrast and correlations & descriptive supported \\
Audited LCZ refinement & 569/569 complete; primary agreement 8.3\% (weak) and 8.3\% (LightGBM) & direct audit comparison \\
CLIMORFA multimodal design & raster, figure-ground, and tabular branches with branch-wise attribution & extensible architecture \\
\bottomrule
\end{tabular}%
}
\begin{minipage}{0.96\textwidth}\vspace{2pt}\scriptsize Note: The summary distinguishes the empirical components of the present analysis from the architecture designed for multimodal extension.\end{minipage}
\end{table}

